\begin{tuluyen}
Để đun sôi hai nồi nước giống nhau, cùng chứa 40 L nước từ nhiệt độ ban đầu 25$^\circ$C, người ta dùng hai bếp: bếp (1) dùng củi, hiệu suất nhiệt 20\%; bếp (2) dùng khí methane, hiệu suất nhiệt 30\%. Biết nhiệt lượng tỏa ra khi đốt cháy 1 gam củi là 20 kJ/g, khi đốt cháy 1 gam methane là 55 kJ/g, nhiệt lượng cần thiết để 1 gam nước lỏng tăng lên 1$^\circ$C là 4,2 J. Tính khối lượng củi và methane cần dùng.
\source{Trích nguyên văn Câu 6.1, đề tuyển sinh chuyên Hóa Vĩnh Long 2026.}
\dapso{\SI{3.15}{\kilo\gram} củi; \(\approx\SI{0.764}{\kilo\gram}\) methane.}
\begin{loigiaiBT}
Coi khối lượng riêng của nước là \SI{1}{\kilo\gram\per\liter}, khối lượng nước trong mỗi nồi là \SI{40}{\kilo\gram}=\SI{40000}{\gram}. Độ tăng nhiệt độ:
\[
\Delta T=100-25=75^\circ\mathrm{C}.
\]
Nhiệt lượng mỗi nồi nước cần nhận:
\[
Q_{\text{cần}}=40000\times4.2\times75=\SI{12600000}{\joule}=\SI{12600}{\kilo\joule}.
\]
Với bếp củi:
\[
Q_{\text{củi}}=\frac{12600}{0.20}=\SI{63000}{\kilo\joule},\qquad
m_{\text{củi}}=\frac{63000}{20}=\SI{3150}{\gram}=\SI{3.15}{\kilo\gram}.
\]
Với bếp methane:
\[
Q_{\ce{CH4}}=\frac{12600}{0.30}=\SI{42000}{\kilo\joule},\qquad
m_{\ce{CH4}}=\frac{42000}{55}\approx\SI{763.6}{\gram}.
\]
\textbf{Vậy,} cần \SI{3.15}{\kilo\gram} củi hoặc khoảng \SI{0.764}{\kilo\gram} methane.
\end{loigiaiBT}
\end{tuluyen}
